\section{A pointwise bound for the global rank cocycles}

% BEGIN D81S1
\paragraph{D81S1. The integral equations for the common global column module.}
Fix a twisted localization and good point at $p=q$ or $d$ as in
I81G1. Write $G=G_E$, let $O$ be the coefficient DVR with uniformizer
$\lambda$ and fraction field $L$, and let $T$ be the given rank
lattice with action $\rho/\chi_1$. Put $V=L\otimes_O T$. Use the
ordered-root cover $X'\to X_i$ and the single integral matrix family
$E_1$ from I81G1. Form the finite direct sum $P$ of one copy of
$R_{x}/\chi_1$ over $L$ for each column at each point $x'$ of the
cover, where $x$ is the image of $x'$. Let $E\in P$ be the vector
of those columns of $E_1$, in that same order, and put
\[
 d_g=gE-E,\qquad Y_0=\sum_{g\in G}O d_g\subset P .
\]
The $d_g$ belong to the common integral coordinate lattice,
so $Y_0$ is finite free over the DVR. It is stable since
$g d_h=d_{gh}-d_g$. For each native generator of $C_\rho$ the map
on the matrix family restricts, after reassembling these columns,
to an equivariant $O$-linear map $h:Y_0\to V/T$ with
$z(g)=h(d_g)$. This is the global column calculation of G79E13:
left multiplication by $R_x(g)/\chi_1(g)$ on a matrix acts on
every column, and on its image the intertwining identity is
precisely the given action on $V/T$. The module $Y_0$ and the
vector $E$ are the same for all generators. Their $O$-linear
combinations have the same presentation.

Set $M=\operatorname{Hom}_O(Y_0,T)$, $M_L=L\otimes_O M$, with
action $(g h)(y)=g h(g^{-1}y)$. Reduction identifies
$M_L/M$ equivariantly with $\operatorname{Hom}_O(Y_0,V/T)$:
lift the images of an integral basis of $Y_0$ to $V$, and observe
that two lifts differ by an element of $M$. In bases of $M$ and
its dual let $D\subset M^\vee$ be the $O$-span of all rows of the
matrices $g-1$ acting on $M$. The module $D$ is Noetherian and
has a generating list drawn from finitely many of those rows,
because each element in a finite generating set uses only finitely
many rows. Stack such rows in a matrix. Write
$\lambda^{v_1},\ldots,\lambda^{v_s}$ for its nonzero Smith
entries, up to units, with exponents $v_a\ge0$.
Integral row and column operations give this diagonal form:
move an entry of least valuation to the first position, use its
divisibility into every remaining entry to clear its row and column,
and repeat on the smaller matrix. In particular $s=0$ is allowed.

A class $u+M$ is fixed if and only if $b(u)\in O$ for every $b\in D$,
while $u\in M_L$ itself is fixed if and only if every $b(u)=0$.
In the Smith coordinates the first $s$ coordinates of a fixed
class range over $\lambda^{-v_a}O/O$ and the others over $L/O$.
Reduction of $M_L^G$ gives exactly the latter coordinates.
Consequently the quotient
\[
 Q= (M_L/M)^G/\operatorname{im}(M_L^G)
   \simeq \bigoplus_{a=1}^s O/\lambda^{v_a}O
\]
has length $\sum_a v_a$, independent of all bases and row choices.
% END D81S1

% BEGIN D81S2
\paragraph{D81S2. The Smith bound before and after simultaneous localization.}
For $h\in(M_L/M)^G$ evaluation $g\mapsto h(d_g)$ is a continuous
cocycle by $d_{gh}=d_g+g d_h$. Indeed $h$ kills a power of $\lambda$
times $Y_0$, which is open in $Y_0$ in its subspace topology:
coordinates in a basis of $Y_0$ extend to continuous linear
coordinates on $P$. The map $g\mapsto gE$ is continuous by the
good-family condition. Evaluation followed by passage to the
cohomology class therefore defines an $O$-linear map to $H^1(G,V/T)$.

Each $R_x$ is semisimple over the specified field $L$ by
$\mathrm{GoodR}$. The character twist preserves stable subspaces
and their complements. A finite direct sum of semisimple modules
is semisimple, so the stable subspace $L Y_0$ has an equivariant
projection $P\to L Y_0$ restricting to the identity there.
Thus every $h_0\in M_L^G=\operatorname{Hom}_{L[G]}(L Y_0,V)$
extends equivariantly to a map $H:P\to V$. The cocycle obtained
from its reduction has the global cobounding vector $H(E)\bmod T$,
since $H(d_g)=gH(E)-H(E)$. It follows that evaluation on cohomology
factors through the finite module $Q$ of D81S1. By I81G1 the
native $C_\rho$ meets the global coboundaries in zero, and D81S1
puts its image in the image of $Q$. Hence for
$e_O=\operatorname{length}_O(O/pO)$ one has
\[
 \frac{\operatorname{length}_O C_\rho}{e_O}
       \le \frac{1}{e_O}\sum_{a=1}^s v_a .
\]
The joint restriction of these same global classes to all
$H^1(E_w,V/T)/H_f^1(E_w,V/T)$ for $w\mid p$ is an image of
$C_\rho$, so its normalized length satisfies the same bound.
This uses the original global matrices and requires no row length
estimate. The integer on the right before normalization belongs to
the chosen point, representation family and projector. Computing
its exact value by the displayed matrix requires a certificate that the chosen
rows generate all the integral equations, as well as a basis of
$Y_0$. For an upper bound it suffices to know the full rational
rank $s$ and to exhibit a nonzero $s$-rowed minor from sampled
equations: in the Smith basis all such minors are divisible by
$\lambda^{\sum_a v_a}$, and integral changes of basis preserve
the ideal of minors. Finite generation proves existence of generating
lists. An uncertified sample of group elements gives no stopping
criterion for having obtained them, and this proof gives no estimate
of the Smith sum in terms of the family parameters.
% END D81S2
